\section{結言}
今回のX観測では、モデル単体の性能競争よりも、それを成立させるインフラと流通経路が前面に出た。NVIDIA・OpenAI・SB Energyの巨大データセンター計画は、生成AIの拡張がGPUだけでなく、電力、長期契約、資金調達まで含む産業基盤の競争になっていることを象徴する。StripeとOpenRouterの買収報道も、モデルを選び、接続し、課金するレイヤー自体が重要な事業領域になったことを示唆している。

同時に、Qwen3.8-27Bのローカル実行やMediaTekの車載対応は、こうした巨大基盤とは逆方向の分散化も進んでいることを示した。クラウド側では計算資源が巨大化する一方、エッジ側では高性能モデルを端末や車両へ持ち込む動きが強まっている。GitHub障害がAIコーディング環境へ波及したことも含め、生成AIの実用性はモデルの能力だけでなく、その周囲のインフラ構成に大きく左右される段階へ入っている。

ウォーターマークや動画生成をめぐる話題も継続しており、能力向上と並行して、生成物の来歴や提供形態をどう扱うかという論点も残る。翌日以降も、巨大AI設備への投資、モデルルーティング事業、オープンモデルのエッジ展開は継続して観測すべき流れになりそうだ。
